% Context from chapters/compression.tex; for mathematical reference, not compilation.
\Qq_T(f)}^2 \leq \norm{f-f_M}^2 + M T^2/4
		\qwhereq
		M = \#\enscond{m}{ \tilde a_m \neq 0}.
	}
\end{prop}

\begin{proof}
The reconstruction rule~\eqref{eq-de-quantization} gives, for $|a_m|\geq T$,
\eq{
	|a_m-\tilde a_m| \leq \frac{T}{2}.
}
By orthogonality,
\eq{
	\norm{f-\Qq_T(f)}^2 = \sum_m (a_m-\tilde a_m)^2 \leq
	\sum_{ |a_m|<T } |a_m|^2 + 
	\sum_{ |a_m| \geq T } \pa{\frac{T}{2}}^2, 
}
The first sum is $\norm{f-f_M}^2$, and the second contains $M$ terms, proving the bound.
\end{proof}


                                        
\subsection{Support Coding}
\label{subsec-support-coding}

To express the error bound~\eqref{eq-error-compression} in terms of a bit budget, we must relate $M$ and $T$ to the cost of encoding the coefficients.

At high compression ratios, only a small number $M$ of the $N$ coefficients remain nonzero. Their support
\eq{
	I_M = \enscond{m}{\tilde a_m \neq 0}
}
can be encoded separately from the nonzero values $q_m \neq 0$ for $m \in I_M$.

The following theorem gives a rate bound for this support-and-value coding strategy.

\begin{thm}\label{thm-compression}
Assume that $f\in L^2$ has best $M$-term approximation error $\norm{f-f_M}^2\leq C M^{-\alpha}$ for some $\alpha>0$. Suppose that, for each sufficiently small $T>0$, all coefficients of magnitude at least $T$ occur among a known first $N(T)$ basis elements, with $N(T)\leq C_0T^{-\gamma}$ for some $\gamma>0$. Then $\Qq_T(f)$ admits a binary description of length $R(T)$ such that, as $R(T)\to\infty$,
\eq{\norm{f-\Qq_T(f)}^2=O\!\left(\left(\frac{\log R(T)}{R(T)}\right)^\alpha\right).}
The decoder is assumed to know the basis and quantization step; describing these parameters adds any required finite header.
\end{thm}

\begin{proof}
\par\smallskip\noindent\textbf{Coefficient and quantization bounds.}\quad
By Proposition~\ref{prop-equiv-comp-decay}, the ordered magnitudes satisfy
\eql{\label{hyp-compression-1}
 d_m\leq C_1m^{-(\alpha+1)/2}.
}
Put $K=T^{-2/(\alpha+1)}$. At most $M=O(K)$ coefficients have magnitude at least $T$. Summing the squared errors of the zero and nonzero quantization cells gives
\eql{\label{eq-compression-error}
 \norm{f-\Qq_T(f)}^2
 \leq\sum_{m\geq1}\min(d_m^2,T^2)
 \leq\lceil K\rceil T^2+\sum_{m>\lceil K\rceil}d_m^2
 =O(K^{-\alpha}).
}
An upper approximation bound is enough for this estimate.

\par\smallskip\noindent\textbf{Finite computation.}\quad
The finite-access assumption is
\eql{\label{eq-compression-lincond}
 \forall m\geq N(T),\qquad |\dotp{f}{\psi_m}|<T.
}
For a bounded signal on the unit cube with periodic boundary conditions, using periodized compactly supported $d$-dimensional wavelets, the coefficient bound is $O(2^{jd/2})$ at fine scales $j\to-\infty$. There are $O(2^{-Jd})$ coefficients at scales coarser than $J$. Retaining scales down to $J$ with $2^{Jd/2}$ proportional to $T$ therefore suffices, and one may choose
\eql{\label{eq-scale-t-n}N(T)=O(T^{-2}).}
More generally, the hypothesis gives
\eql{\label{hyp-compression-2}
 N(T)=O(K^{\gamma(\alpha+1)/2}).
}
The ordering here is a fixed computable ordering of basis elements, not the unknown ordering by coefficient magnitude.

\par\smallskip\noindent\textbf{Support and value coding.}\quad
Encode each of the $M$ nonzero indices with $\lceil\log_2N\rceil$ bits:
\eql{\label{eq-nbits-1}
 R_{\mathrm{ind}}=M\lceil\log_2N\rceil=O(K\log K).
}
Since $|a_m|\leq\norm{f}_2$, every quantized value lies in $[-A,A]\cap\ZZ$ with $A\leq\norm{f}_2/T$. Thus
\eql{\label{eq-nbits-2}
 R_{\mathrm{val}}\leq M\lceil\log_2(2A+1)\rceil=O(K\log K).
}
A header encoding $M$, $N$, and the integer amplitude bound $A$ has logarithmic size and is absorbed in this bound. Consequently,
\eql{\label{eq-bound-bits}R=O(K\log K).}
Because the inverse of $x\mapsto x\log x$ is asymptotic to $r/\log r$, this implies
\eql{\label{eq-compression-bits}K\geq cR/\log R}
for large $R$. Combining this with~\eqref{eq-compression-error} proves the claim.
\end{proof}


The theorem links compression rates to nonlinear approximation rates and motivates choosing a basis adapted to the signal model $\Th$.
 
A practical compression algorithm operates on a discrete signal or image of size $N$. It computes $N$ transform coefficients $\{\dotp{f}{\psi_m}\}_{0\leq m<N}$ from these data. Section~\ref{sec-fast-wavelet-transform} describes the discrete wavelet transform and gives a compatibility condition~\eqref{eq-hyp-wav-samples} under which these discrete coefficients coincide with the continuous inner products.


\myfigure{
\tabquatre{
\image{approximation}{.24}{compression-original}&
\image{approximation}{.24}{compression-original-zoom}&
\image{approximation}{.24}{compression-support}&
\image{approximation}{.24}{compression-result}\\
Original $f$ & Zoom & Wavelet support & Compressed $\Qq_T(f)$
}
}{ 
	Image compression using wavelet support coding.  
}{fig-compression}



                                                                 
\section{Entropy Coding}
\label{subsec-entropic-coding}

Entropy coding exploits the uneven distribution of the quantized coefficients to reduce the expected file size: zeros are frequent, while large values are rare. Shannon's coding theory \cite{shannon-information} quantifies the resulting savings.

We refer to Section~\ref{sec-shannon-source} for the theoretical foundations of the methods described below.


  
\paragraph{Probabilistic modeling.}

The quantized coefficients $q_m \in \{-A,\ldots,A\}$ take values in an alphabet of $Q=2A+1$ elements.
A code maps the coefficient sequence to a binary string:
\eq{
	\{q_m\}_m \longmapsto \{0,1,1,\ldots,0,1\} \in \{0,1\}^R.
}
We first model the $q_m$ as independent draws from a known probability distribution:
\eq{
	\PP(q_m=i)=p_i \in [0,1].
}

  
\paragraph{Huffman code.}

A Huffman code assigns a variable-length binary string to each symbol:
\eq{
	q_m=i \in \{-A,\ldots,A\} \longmapsto c_i \in \{0,1\}^{|c_i|}
}
where $|c_i|$ is the codeword length. Codeword lengths are assigned so that a less probable symbol receives no shorter a codeword than a more probable symbol.
Huffman coding minimizes the expected length among binary prefix codes. Shannon's construction gives lengths $\lceil-\log_2p_i\rceil$ for positive-probability symbols, so optimality of Huffman's code implies
\eq{\mathbb E R=N\sum_i p_i|c_i|\leq N(\Ee(p)+1),}
where $\Ee$ is the entropy of the distribution, defined as
\eq{
	\Ee(p) = -\sum_i p_i \log_2(p_i).
}
Figure~\ref{fig-entropy} illustrates how entropy depends on the distribution. Concentrated distributions have low entropy, as is often the case for quantized wavelet coefficients because many are zero.

Symbol-by-symbol Huffman coding uses an integer number of bits per symbol. This constraint can cause a substantial gap above entropy when one symbol is very probable. Arithmetic coding encodes whole sequences and can achieve an expected 